\chapter{Donnees et Pipeline}

\section{Vue d'ensemble}
Le pipeline retenu pour les resultats finaux suit les etapes suivantes:
\begin{enumerate}[leftmargin=2em]
  \item collecte Wikipedia (\texttt{scripts/wiki\_fetch.py}),
  \item chunking des textes longs (\texttt{scripts/chunk\_texts.py}),
  \item extraction timeline par LLM (\texttt{scripts/extract\_timelines.py}),
  \item traduction FR->EN des timelines FR, sauvegardees en \texttt{*.fr\_en.timeline.json}
        (\texttt{scripts/translate\_fr\_timelines\_to\_en.py}),
  \item construction des paires \texttt{en <-> fr\_en}
        (\texttt{scripts/make\_pairs.py}, \texttt{scripts/build\_pairs\_autotrain.py}),
  \item fine-tuning et evaluation similarite (\texttt{scripts/eval\_similarity.py}),
  \item generation controlee (\texttt{scripts/generate\_analogues.py}) et UI (\texttt{app.py}).
\end{enumerate}

Une branche complementaire \texttt{merge -> final\_en} existe egalement
(\texttt{scripts/merge\_multilingual\_timelines.py},
\texttt{scripts/finalize\_monolingual\_en\_timelines.py}),
mais elle n'est pas la base des scores finaux reportes.

\section{Statistiques principales (pipeline retenu)}
\begin{table}[H]
  \centering
  \caption{Statistiques de la branche \texttt{en/fr\_en} retenue}
  \begin{tabular}{lr}
    \toprule
    Indicateur & Valeur \\
    \midrule
    Timelines \texttt{en} disponibles & 249 \\
    Timelines \texttt{fr\_en} disponibles & 249 \\
    Entites retenues pour paires \texttt{en <-> fr\_en} & 248 \\
    Paires positives & 248 \\
    Paires negatives & 744 \\
    Total paires & 992 \\
    \bottomrule
  \end{tabular}
\end{table}

Les statistiques proviennent de:
\texttt{data/datasets/fr\_en\_timelines.stats.json} et
\texttt{data/datasets/pairs\_en\_fr\_en\_mvp.stats.json}.

\section{Gestion du contexte long}
Le chunking decoupe les biographies en segments coherents, avec nettoyage des
sections peu utiles (references, liens externes, etc.). Cette strategie reduit
les erreurs de parsing et ameliore la couverture factuelle sur les textes longs.
